\documentclass[11pt]{article}

\usepackage[margin=1in]{geometry}
\usepackage{amsmath,amssymb}
\usepackage{longtable}
\usepackage{array}

\newcommand{\toprule}{\hline}
\newcommand{\midrule}{\hline}
\newcommand{\bottomrule}{\hline}
\newcommand{\href}[2]{\texttt{\detokenize{#2}}}
\newcommand{\url}[1]{\texttt{\detokenize{#1}}}

\title{Overlap Guided Adaptive Fractionation}
\author{
  Yoel P\'erez Haas\thanks{Shared first author.} \and
  Lena Kretzschmar\footnotemark[1] \and
  Bertrand Pouymayou \and
  Stephanie Tanadini-Lang \and
  Jan Unkelbach
}
\date{
  Department of Radiation Oncology, University Hospital of Zurich, Zurich, Switzerland\\
  \vspace{0.5em}
  November 5, 2025
}

\begin{document}
\maketitle

\begin{center}
\textbf{Internal LaTeX transcription for repository reference}
\end{center}

\noindent\textit{Editorial note.} This file is a manuscript-style transcription of the 2025 paper assets in this repository. Numbered equations were transcribed from the equation screenshots in [\texttt{docs/papers/2025/eqtns}] and cross-checked against the extracted manuscript text. Equation~\eqref{eq:sigma-map-overlap} is preserved as written in the manuscript assets, even though its notation is not fully consistent with equation~\eqref{eq:gamma-prior-overlap}.

\begin{abstract}
\textbf{Purpose.}
Online-adaptive, magnetic-resonance-guided radiotherapy on hybrid MR-linear accelerators enables stereotactic body radiotherapy (SBRT) of abdominal and pelvic tumors with large interfractional motion. However, overlaps between planning target volume (PTV) and dose-limiting organs at risk (OARs) often force compromises in PTV coverage. Overlap-guided adaptive fractionation (AF) leverages daily variations in PTV/OAR overlap to improve PTV coverage by administering variable fraction doses based on measured overlap volume. This study aims to assess the potential benefits of overlap-guided AF.

\textbf{Methods.}
We analyzed 58 patients with abdominal or pelvic tumors having received 5-fraction MR-guided SBRT ($> 6$ Gy/fraction), in whom PTV overlap with at least one dose-limiting OAR (bowel, duodenum, stomach) occurred in at least one fraction. Dose-limiting OARs were constrained to $1\mathrm{cc} \leq 6$ Gy per fraction, rendering overlapping PTV volumes underdosed. AF aims to reduce this underdosage by delivering higher doses to the PTV on days with less overlap volume, and lower doses on days with more. PTV coverage gain compared to uniform fractionation was quantified by the area above the PTV dose-volume-histogram curve and expressed in ccGy ($1\,\mathrm{ccGy}=1\,\mathrm{cc}$ receiving $1\,\mathrm{Gy}$ more). The optimal dose for each fraction was determined through dynamic programming by formulating AF as a Markov decision process.

\textbf{Results.}
PTV/OAR overlap volume variation (standard deviation) varied substantially between patients (0.02--5.76 cc). Algorithm-based calculations showed that 55 of 58 patients benefited in PTV coverage from AF. Mean cohort benefit was 2.93 ccGy (range $-4.44$ to 22.42 ccGy). Higher PTV/OAR overlap variation correlated with larger AF benefit.

\textbf{Conclusion.}
Overlap-guided AF for abdominal and pelvic SBRT is a promising strategy to improve PTV coverage without compromising OAR sparing. Since the benefit of AF depends on PTV/OAR overlap variation, which is low in many patients, the mean cohort advantage is modest. However, well-selected patients with marked PTV/OAR overlap variation derive a relevant dosimetric benefit. Prospective studies are needed to evaluate AF feasibility and quantify clinical benefits.
\end{abstract}

\section{Introduction}

The recent technological development of a hybrid MR-linear accelerator (MR-linac), i.e. a combination of an MR scanner, a 6 MV linear accelerator, and a multileaf collimator for intensity modulation, makes real-time MR-guided radiotherapy feasible \cite{ref1,ref2,ref3,ref4}, and enables the delivery of high-precision stereotactic body radiotherapy (SBRT) to traditionally more challenging target volumes, e.g. mobile abdominal lesions such as pancreatic cancer, as well as liver or adrenal metastases \cite{ref5,ref6,ref7,ref8}. A main advantage of MR-guided radiotherapy is the visualization of patient anatomy with superior soft-tissue contrast compared to CT-guided radiotherapy. This enables daily, pre-treatment adaptation of gross tumor volume (GTV) and organs at risk (OAR).

A challenge in the application of MR-guided SBRT for abdominal and pelvic target volumes is posed by closeness or overlap of the planning target volume (PTV) with adjacent, highly vulnerable OARs, e.g. bowel and stomach. This issue, if it arises, cannot be improved by simply adjusting the treatment plan to the geometry of the day, as is routinely performed in MR-guided radiotherapy \cite{ref9}. Usually, compromises in PTV dose coverage, or a reduction of prescription dose, are necessary, with both options being undesirable due to their potentially negative effects on local tumor control.

However, owing to interfraction mobility of both GTV and OARs, the geometry between them does not remain static over the course of treatment, as demonstrated in Figure~\ref{fig:overview-patient3}. Adaptive fractionation (AF) \cite{ref10,ref11,ref12} is one approach to utilize these daily geometric changes to provide better dose coverage of the PTV without compromising sparing of OARs. Its general concept is to apply more dose to the PTV on days with more favorable PTV/OAR geometry, and less dose on days with less favorable PTV/OAR geometry.

AF has previously been explored by our group in the context of MR-guided radiotherapy \cite{ref13}. In this prior work, day-to-day variation of the geometry was quantified via a so-called sparing factor defined as the quotient of the dose received by the dose-limiting OAR ($D_{1\mathrm{cc}}$) and the dose received by the PTV ($D_{95}$) in each fraction. AF was implemented by upscaling the dose distribution for a small sparing factor and downscaling for a large sparing factor. The sparing factor is smaller the larger the distance between the PTV and the OAR. However, the sparing factor is approximately 1 if the PTV and the OAR overlap and the minimum dose in the PTV equals the maximum dose in the OAR. This limits the applicability of AF for patients in whom PTV and dose-limiting OAR have direct volume overlap in each fraction, because there is little variation in the sparing factor even if the overlap volume varies substantially.

Based on these previous observations, we developed an alternative AF approach for cases in which direct overlaps between PTV and dose-limiting OARs are present, by investigating overlap volume as the essential descriptor of geometric variation and the main leverage point for AF. We demonstrate this new approach using data from patients who previously received 5-fraction SBRT in the form of MR-guided radiotherapy at the MR-linac in our department for abdominal and pelvic target volumes, and whose PTV showed overlap with at least one dose-limiting hollow OAR: bowel, duodenum, or stomach.

\begin{figure}[ht]
\centering
\fbox{\parbox{0.8\linewidth}{Figure omitted in this text transcription. See the original PDF for the rendered figure.}}
\caption{MR-guided SBRT examples for patient 3. Panel (a) shows fraction 2/5 and panel (b) shows fraction 5/5. The blue contour represents the GTV, the red contour the PTV, the green contour the dose-limiting OAR, and the yellow shading the overlap region.}
\label{fig:overview-patient3}
\end{figure}

\section{Patient Selection and Treatment Plans}

The data for this analysis stems from all patients with abdominal and pelvic primary or secondary tumors who received SBRT at the MR-linac (MRIdian, ViewRay) in our department between 2019 and 2024. Patients were included in the analysis if the PTV was located in direct proximity to the dose-limiting OARs bowel, stomach, and/or duodenum, overlapping with at least one of these organs in at least one fraction of treatment, and thus necessitating a PTV coverage compromise to respect OAR constraints. Additionally, to qualify for inclusion, SBRT had to have been administered in 5 fractions with a minimum dose of 6 Gy per fraction, regardless of the prescribed isodose level. To limit complexity, treatment courses were excluded if more than one lesion was irradiated in the same plan.

We isolated a total of 58 treatment courses that met these criteria and extracted volumetric and dosimetric information on GTV, PTV, and relevant OARs for a total of 6 treatment plans per course (the initial plan based on the planning MRI, plus the 5 delivered plans) from the treatment planning system. Treatment planning and delivery for all included patients followed institutional protocols: each patient underwent either MR-only or combined MR and CT simulation, as well as daily MR imaging for online adaptive radiotherapy. Tumors and OARs within a 2 cm margin around the GTV were recontoured prior to each fraction according to institutional guidelines. Adaptive treatment plans were generated daily to account for interfractional anatomical changes, with prescription doses remaining constant across all fractions.

Overlap between PTV and the relevant dose-limiting OARs was calculated by creating an overlap structure for each overlapping OAR with the PTV and measuring its volume. If more than one dose-limiting OAR showed overlap, all separate overlap volumes were added together to create a total overlap volume. This was possible because bowel, duodenum, and stomach share the same OAR objectives in the institutional protocol. All considerations on AF below were made based on this total overlap volume.

The most common malignancies in the cohort were primary and secondary tumors of the pancreas ($n=24$), followed by adrenal metastases ($n=13$) and abdominal or pelvic lymph node metastases ($n=13$); the remaining patients had abdominal tumor manifestations that were not more closely defined ($n=8$), e.g. non-lymph-node-associated soft-tissue growths. Thirty-two patients had 35 Gy prescribed to the PTV as total dose, 14 received 33 Gy, 11 received 40 Gy, and one patient was prescribed 45 Gy. Due to the daily-adaptive nature of the workflow, PTV volumes showed fluctuations over the course of radiotherapy, depending on the daily size of the GTV, and influenced by other factors such as interobserver variability between contouring physicians and the image quality of the daily MRI scan. PTV volume interfraction variation, i.e. standard deviation, was generally less than 10\% of PTV volume; only 7 patients showed a variation of PTV size greater than 10\% over the course of SBRT, with a maximum of 18\% PTV volume in a single patient due to tumor swelling and reduced visibility conditions creating uncertainties in the adaptive planning MRIs.

Mean PTV volumes per course varied and ranged from 1.51 cc to 286 cc in the patient cohort. No clear dependency between PTV volume and tumor entities or locations was found.

\subsection{Overlap Variation}

PTV/OAR overlap volume variation, i.e. its standard deviation, showed high variability between patients, ranging from 0.02 cc up to 5.76 cc. Figure~\ref{fig:top-overlap-variation} shows the ten patients with the largest overlap variations in the cohort and their respective PTV/OAR overlaps, stratified by fraction; a detailed chart of the overlaps of all patients can be found in the appendix.

\begin{figure}[ht]
\centering
\fbox{\parbox{0.8\linewidth}{Figure omitted in this text transcription. See the original PDF for the rendered figure.}}
\caption{Variation in overlap volumes for the ten patients with the highest standard deviation, sorted by magnitude of variation. Each point represents the total calculated overlap volume for one fraction, including the initial simulation. The red line represents the mean overlap volume throughout each treatment course.}
\label{fig:top-overlap-variation}
\end{figure}

Figure~\ref{fig:gamma-prior-2025} shows the distribution of the patient-specific standard deviations for all 58 patients. A gamma distribution was fitted to this empirical data, which later serves as a population-level prior in the AF framework introduced in section~3.2.

\begin{figure}[ht]
\centering
\fbox{\parbox{0.8\linewidth}{Figure omitted in this text transcription. See the original PDF for the rendered figure.}}
\caption{Gamma prior fitted to the observed overlap volume standard deviations of all patients, resulting in shape parameter $k = 1.07$ and scale parameter $\theta = 0.78$.}
\label{fig:gamma-prior-2025}
\end{figure}

Both in absolute and in relative numbers (calculated as part of PTV volume) we found a statistically significant correlation between mean PTV/OAR overlap over the course of treatment and the variation (standard deviation) of PTV/OAR overlap, as illustrated in Figure~\ref{fig:overlap-correlation}. Overlap variation did not show a relevant dependency on mean PTV volume or tumor entity.

\begin{figure}[ht]
\centering
\fbox{\parbox{0.8\linewidth}{Figure omitted in this text transcription. See the original PDF for the rendered figure.}}
\caption{Correlation between mean PTV/OAR overlap and overlap variation, shown in both absolute and relative form. The colors represent tumor locations.}
\label{fig:overlap-correlation}
\end{figure}

Looking at the relative PTV/OAR overlaps and their respective standard deviations, we found an interesting trend of small target volumes, especially lymph nodes, having a high relative overlap variation for their size, despite absolute overlap volumes and variations being generally small in these patients. This led us to further investigate the potential benefits of AF not only in absolute but also in relative numbers.

\section{Adaptive Fractionation}

AF aims to adjust the dose delivered in each treatment fraction based on daily anatomical variations, in this case the overlap between the PTV and OARs. The key idea is to escalate the dose when the anatomy is favorable (small overlap) and reduce the dose when there is significant overlap.

Treatment planning is subject to the following clinical objectives and constraints:

\begin{enumerate}
\item The dose in the OAR, including the overlap volume, is constrained to 6 Gy in each fraction. In this study, we consider hypofractionated treatments consisting of $F=5$ fractions. Thus, this corresponds to a commonly used cumulative dose constraint of 30 Gy in 5 fractions.
\item The total cumulative dose prescribed to the PTV is larger than 30 Gy but the OAR constraint is prioritized over PTV coverage. Consequently, the PTV is underdosed in the overlap region.
\item The per-fraction dose prescribed to the PTV, $d_t$, must lie within the range $d_t \in [d_{\min}, d_{\max}]$, where $d_{\min}$ is set equal to the OAR constraint of 6 Gy.
\end{enumerate}

The goal is to determine the daily prescription doses $d_t$ for each fraction $t$ such that the overall PTV underdosage is minimized. To that end, the following quantitative measure of PTV underdosage is introduced.

\subsection{Measure of Underdosage}

To characterize interfractional anatomical variability, we quantify the volumetric overlap between the PTV and OARs. Specifically, let $o_t$ denote the PTV volume overlapping with the relevant OAR on treatment day $t$, measured in cc. As an approximate measure for the loss in PTV coverage in fraction $t$, we introduce a cost function $c_t$ defined as

\begin{equation}
c_t = o_t \cdot (d_t - d_{\min}).
\label{eq:underdosage-cost}
\end{equation}

Here, $d_t$ is the prescribed dose to the PTV in fraction $t$ and $d_{\min}$ is the dose that is deliverable to the overlapping region without violating the per-fraction OAR constraint. The cost $c_t$ is thus proportional to the amount by which the prescribed fraction dose exceeds $d_{\min}$.

A schematic representation of this concept is shown in Figure~\ref{fig:dvh-schematic}. The illustration shows the dose-volume histogram (DVH) of the PTV for a prescription dose $d_t > d_{\min}$. The green area above the DVH curve represents the underdosage of the PTV due to the dose constraint in the overlap volume. The red box represents the approximation of the area above the DVH curve as defined by the metric $c_t$. Thus, the cost is measured in units of ccGy, where an improvement of 1 ccGy indicates that 1 cc of the PTV volume receives 1 Gy more. The green area can only be determined after a treatment plan has been created for a prescription dose $d_t$, whereas the red area represents an approximation that can be estimated based on the overlap volume $o_t$ alone prior to treatment planning.

\begin{figure}[ht]
\centering
\fbox{\parbox{0.8\linewidth}{Figure omitted in this text transcription. See the original PDF for the rendered figure.}}
\caption{Schematic illustration of the underdosage measure and comparison of DVHs for different dose prescriptions and overlap scenarios.}
\label{fig:dvh-schematic}
\end{figure}

The overall quality of an adaptive fractionation scheme can be evaluated by aggregating the daily costs over the entire treatment course. We define the total accumulated cost as

\begin{equation}
c_{\mathrm{total}} = \sum_{t=1}^{F} o_t (d_t - d_{\min}).
\label{eq:total-cost}
\end{equation}

To further illustrate the rationale for overlap-based AF, consider two fractions with different overlap volumes. Suppose $o_1 = 10$ cc in the first fraction and $o_2 = 3$ cc in the second fraction. Let us further assume a uniform fractionation scheme prescribes a constant dose of $d_{\mathrm{ref}} = 8$ Gy in both fractions, with an OAR constraint of $d_{\min}=6$ Gy. The total cost over two fractions for uniform fractionation is

\[
\text{Uniform Fractionation:}\qquad d_1 = d_2 = 8~\mathrm{Gy}, \qquad
c_{\mathrm{total}} = 10\,\mathrm{cc}\cdot(8-6)\,\mathrm{Gy} + 3\,\mathrm{cc}\cdot(8-6)\,\mathrm{Gy} = 26\,\mathrm{ccGy}.
\]

In contrast, an adaptive strategy could deliver 6 Gy in the first unfavorable fraction and 10 Gy in the second favorable fraction. Since the dose of 6 Gy respects the OAR constraint, no underdosage occurs in fraction one regardless of the overlap volume. The total cost over the two fractions for adaptive fractionation is

\[
\text{Adaptive Fractionation:}\qquad d_1 = 6~\mathrm{Gy},\quad d_2 = 10~\mathrm{Gy}, \qquad
c_{\mathrm{total}} = 10\,\mathrm{cc}\cdot(6-6)\,\mathrm{Gy} + 3\,\mathrm{cc}\cdot(10-6)\,\mathrm{Gy} = 12\,\mathrm{ccGy}.
\]

In this simplified example, adaptive fractionation reduces the cost by 14 ccGy compared to uniform fractionation. Even though the underdosage in fraction two increases when 10 Gy instead of 8 Gy is prescribed, this is overcompensated for by avoiding the substantial underdosage in fraction one that would otherwise occur due to the large overlap volume.

\subsection{Probability Distribution for Overlap Volumes}

The central challenge in AF is that future anatomical configurations are not known in advance. Adaptive fractionation relies on anticipated overlap volumes for future fractions when determining the optimal dose in the current fraction. In this work, we assume that the daily overlaps follow a normal distribution $o_t \sim \mathcal{N}(\mu,\sigma^2)$, which is updated iteratively during treatment. As in \cite{ref13}, we estimate the patient-specific parameters $\mu$ and $\sigma$ based on observed overlap values up to the current fraction.

Because mean overlap volumes can differ significantly across patients, we use the maximum likelihood estimator to compute the patient-specific mean, updating it with each observed value:

\begin{equation}
\mu_t = \frac{1}{t+1}\sum_{\tau=0}^{t} o_{\tau}.
\label{eq:mu-overlap}
\end{equation}

Here, $o_0$ denotes the overlap volume observed in the planning MR.

The standard deviation estimated from the observed overlap volumes up to fraction $t$ is given by

\begin{equation}
\sigma_t^{\mathrm{pat}} = \sqrt{\frac{1}{t+1}\sum_{\tau=0}^{t}(o_{\tau} - \mu_t)^2}.
\label{eq:sigma-pat-overlap}
\end{equation}

Since the standard deviation strongly influences the adaptive fractionation policy, we adopt a Bayesian approach to regularize early estimates. Specifically, we use a population-based gamma prior for $\sigma$:

\begin{equation}
f(\sigma; k, \theta) = \frac{1}{\Gamma(k)\theta^k}\sigma^{k-1}\exp\left(-\frac{\sigma}{\theta}\right).
\label{eq:gamma-prior-overlap}
\end{equation}

The hyperparameters $k$ (shape) and $\theta$ (scale) were estimated empirically from the distribution of standard deviations across the retrospective cohort of 58 patients (Figure~\ref{fig:gamma-prior-2025}), yielding $k=1.07$ and $\theta=0.78$.

Given this prior and the observed overlaps up to fraction $t$, we compute the maximum a posteriori estimate of $\sigma_t$. The likelihood is based on a normal model with unknown variance, and the resulting maximum a posteriori estimate can be obtained numerically by maximizing the posterior:

\begin{equation}
\sigma_t = \arg\max_{\sigma}\left[\frac{\sigma^{k-1}}{\sigma^{t-1}} \exp(-\theta \sigma)\exp\left(\frac{(\sigma_t^{\mathrm{pat}})^2}{2\sigma_t^2}\right)\right].
\label{eq:sigma-map-overlap}
\end{equation}

The a posteriori estimator balances patient-specific observations with population-informed regularization, yielding more robust estimates, especially in the early treatment phase when few overlap volumes were observed.

\subsection{Markov Decision Process Model}

To solve the adaptive fractionation problem under uncertainty about future overlaps, we model the decision process as a Markov decision process (MDP) similarly to \cite{ref13}. The goal is to find a dose selection policy that minimizes the expected cumulative cost over $F$ treatment fractions, while delivering the prescribed cumulative dose to the PTV.

The MDP is described by states, actions, state transitions, and cost:

\paragraph{State.}
In each fraction $t$, the state is defined as the tuple $s_t = (o_t, D_t)$ where
\begin{itemize}
\item $o_t$ is the observed PTV/OAR overlap volume in fraction $t$,
\item $D_t = \sum_{\tau=1}^{t-1} d_{\tau}$ is the cumulative dose delivered to the PTV up to, but not including, fraction $t$.
\end{itemize}
This state encapsulates the current geometric condition and treatment history, both of which influence the optimal dose decision.

\paragraph{Action.}
The action is the dose $d_t$ prescribed to the PTV in fraction $t$. The action space is bounded:

\begin{equation}
d_t \in [d_{\min}, d_{\max}].
\label{eq:action-space-overlap}
\end{equation}

In this work, the lower bound $d_{\min}$ is set to 6 Gy, equivalent to the dose constraints in the OAR, and the upper bound is 10 Gy.

\paragraph{State transition.}
After applying dose $d_t$ in state $s_t = (o_t, D_t)$, the state transitions to a new state $s_{t+1} = (o_{t+1}, D_{t+1})$, where $o_{t+1} \sim \mathcal{N}(\mu_t,\sigma_t^2)$ is drawn from the patient-specific normal distribution and $D_{t+1} = D_t + d_t$ is the updated cumulative PTV dose. Since $o_{t+1}$ is stochastic, the transition is probabilistic.

\paragraph{Cost.}
The immediate cost in each fraction is defined by the cost introduced above:

\begin{equation}
c_t = o_t \cdot (d_t - d_{\min}).
\label{eq:immediate-cost-overlap}
\end{equation}

In addition, we impose a terminal cost $c_{F+1}$ to ensure that the total prescribed PTV dose $D_{\mathrm{pres}}$ is delivered exactly:

\begin{equation}
c_{F+1} =
\begin{cases}
0 & \text{if } D_F + d_F \geq D_{\mathrm{pres}}, \\
\infty & \text{otherwise.}
\end{cases}
\label{eq:terminal-cost-overlap}
\end{equation}

\subsection{Dynamic Programming Algorithm}

The optimal policy $\pi_t(o,D)$ assigns to each state $(o,D)$ the optimal action that minimizes the expected cumulative cost until the end of treatment in fraction $F$. To compute the optimal policy for adaptive fractionation, we apply dynamic programming (DP), which recursively computes the optimal dose to deliver in each fraction by minimizing the sum of immediate cost and expected future cost \cite{ref14}. DP is based on the value function $v_t(o,D)$, which quantifies the expected future cost when starting from state $(o,D)$ in fraction $t$ and following the optimal policy from there on. Optimal policy and value function can then be calculated in a backward recursion starting in the last fraction $t=F$ and iterating backwards to $t=1$. One iteration of DP is given by

\begin{equation}
v_t(o,D) = \min_{d \in [d_{\min}, d_{\max}]}\left[o\cdot(d-d_{\min}) + \sum_{o'} P(o') v_{t+1}(o', D+d)\right],
\label{eq:value-overlap}
\end{equation}

\begin{equation}
\pi_t(o,D) = \arg\min_{d \in [d_{\min}, d_{\max}]}\left[o\cdot(d-d_{\min}) + \sum_{o'} P(o') v_{t+1}(o', D+d)\right].
\label{eq:policy-overlap}
\end{equation}

The value function is initialized through the terminal cost, i.e. $v_{F+1}(o,D)=c_{F+1}(o,D)$. Here, $P(o')$ denotes the probability distribution over the unknown overlap volume in the next fraction, which is given by the normal distribution $\mathcal{N}(\mu,\sigma^2)$. The sum over $o'$ accounts for the stochastic nature of future anatomy.

To compute $v_t$ and $\pi_t$ in practice, both the state space, i.e. $D$ and $o$, and the action space $[d_{\min}, d_{\max}]$ are discretized. Note that we make the approximation that the estimates of $\mu$ and $\sigma$ are not part of the state of the MDP but are assumed to be constant over future fractions. To apply AF to a given patient with a given sequence of overlap volumes, we update $\mu$ and $\sigma$ in each fraction as described in section~3.2 and recompute the optimal policy.

\subsection{Benefit Evaluation}

To assess the benefit of AF, we first evaluate fractionation schemes based on the cumulative cost $c_{\mathrm{total}}$ as a measure of PTV underdosage. In a second step, replanning of selected patients is performed.

\paragraph{1. Benchmarking based on cost.}
We compare $c_{\mathrm{total}}$ for AF to two benchmarks:
\begin{itemize}
\item \textbf{Uniformly fractionated treatment.} This corresponds to standard treatment as delivered in practice, where the dose is uniformly fractionated across all days. The total accumulated cost for this treatment serves as a baseline for comparison.
\item \textbf{Upper bound.} This represents an idealized scenario in which the full sequence of overlap volumes $o_t$ is known in advance. Based on this, we compute the optimal set of fraction doses that minimize the total cost. This plan serves as an upper bound on what any realistic adaptive policy could achieve.
\end{itemize}

\paragraph{2. Evaluation through replanning.}
In a second step, we assess the benefit of AF based on the actually achievable PTV DVHs. To that end, replanning of all fractions is performed in the MR-linac treatment planning system for selected patients, using the prescribed doses per fraction recommended by the AF algorithm. For each replanned fraction, we extract the DVH and compute the underdosage in the PTV via the area above the DVH curve. A complete treatment course is characterized by summing the areas above the DVH across all fractions. By comparing these replanned AF-based fractions with the clinically delivered ones, we can quantify the actually realized dosimetric benefit of AF.

\section{Illustration of Adaptive Fractionation}

To illustrate the behavior of the AF algorithm, we present two patients in detail for whom we characterize the optimal policy and discuss the AF scheme and its benefit.

\subsection{Patient 3 (33 Gy Prescription)}

Patient 3 was prescribed a total dose of 33 Gy in 5 fractions. Given that 30 Gy are delivered by applying the minimum dose in each fraction, the prescribed PTV dose exceeds this minimum dose by only 3 Gy. Consequently, the maximum dose that can be delivered in any fraction is 9 Gy. The PTV/OAR overlap volumes were $o_0 = 9.08$ cc on the planning MRI, followed by $[19.79, 6.02, 9.45, 19.59, 12.62]$ cc across the five treatment fractions.

At the start of the first fraction, the estimated mean and standard deviation of the probability distribution over future overlap volumes are estimated from the planning scan and the first fraction scan. This results in $\mu_1 = 14.4$ cc and $\sigma_1 = 3.2$ cc, respectively, as computed via equation~\eqref{eq:sigma-map-overlap}. The resulting overlap distribution is shown in Figure~\ref{fig:patient3-policy}.

Figure~\ref{fig:patient3-policy} shows the optimal policy and value function for patient 3. In the first fraction, since no dose has yet been delivered, the full range of allowed doses (6--9 Gy) is available. The algorithm yields a binary policy: the maximum dose (9 Gy) is prescribed if the overlap is less than 11.7 cc, and the minimum dose (6 Gy) is prescribed if the overlap exceeds 11.8 cc. This threshold lies at a cumulative probability of roughly 20\%, i.e. the dose of 9 Gy is delivered if the probability of observing a lower overlap volume in the future is less than 20\%.

The binary nature of the optimal policy originates from the immediate cost being a linear function of the prescribed dose. As a consequence, the optimal dose to deliver is either the minimum dose or the maximum dose. The corresponding value function in fraction one is constant for overlaps above 11.7 cc where 6 Gy is delivered in the first fraction. For overlaps below 11.7 cc where 9 Gy is delivered in the first fraction, the value function is linear in the overlap volume.

In the second fraction, a favorable overlap of $o_2 = 6.02$ cc is observed. The updated distribution has a lower mean ($\mu_2 = 11.63$ cc) and a slightly higher standard deviation ($\sigma_2 = 3.88$ cc). Again, the policy exhibits a sharp decision boundary. For overlap volumes larger than 9.2 cc, the algorithm suggests delivering the minimum dose of 6 Gy irrespective of the dose delivered in the first fraction. For overlap volumes smaller than 9.2 cc, corresponding to the 27th percentile, the policy depends on the previously accumulated dose in the PTV: if only 6 Gy was delivered previously, the algorithm chooses to deliver 9 Gy now. However, if 9 Gy had already been delivered earlier, the policy prescribes only 6 Gy to avoid exceeding the cumulative PTV dose prescription.

The optimal policy in fractions three and four is qualitatively analogous to fraction two. The policy exhibits a sharp decision boundary. For overlap volumes above a threshold, the minimum dose of 6 Gy is delivered; for overlap volumes below the threshold, the maximum dose needed to reach the total prescribed PTV dose is delivered. However, the threshold shifts to larger percentiles of the overlap distribution for later fractions. In fraction four, the decision boundary lies at the 50th percentile, which is intuitive. If there is a probability larger than 50\% that the observed overlap in the last fraction is lower than the overlap observed in the fourth fraction, 6 Gy will be delivered in the fourth fraction. Otherwise, 9 Gy will be delivered.

Finally, in the fifth fraction, no further adaptation is possible. The remaining dose required to meet the total PTV prescription is delivered, regardless of the observed overlap.

\begin{figure}[ht]
\centering
\fbox{\parbox{0.8\linewidth}{Figure omitted in this text transcription. See the original PDF for the rendered figure.}}
\caption{Optimal policies and value functions for patient 3. Panels show the optimal policies for fractions one through five and the expected cost in fraction one. The black crosses mark the observed overlap and the resulting dose decision.}
\label{fig:patient3-policy}
\end{figure}

Table~\ref{tab:patient3-2025} summarizes the results of the AF treatment for patient 3, alongside two benchmarks: the upper bound and the uniformly fractionated reference treatment. The upper bound corresponds to a treatment that delivers 9 Gy in the fraction with the smallest overlap. Although this is not known in advance, the AF algorithm takes the right decision for this patient. For the sequence of overlap volumes observed for this patient, the AF algorithm identifies the overlap volume in fraction two as small compared to what is expected based on the planning scan and fraction one, and delivers the maximum dose of 9 Gy accordingly. Since indeed no smaller overlap volumes are observed in fractions three to five, AF achieves the lowest possible PTV underdosage given by the upper bound.

In contrast, the reference treatment follows uniform fractionation and applies the same dose (6.6 Gy) in every fraction, regardless of daily anatomical variations. As a result, the total cost accumulated is 40.5 ccGy, compared to just 18 ccGy for the AF plan. This translates to a reduction in cost of 22.5 ccGy, representing better coverage of the PTV.

\begin{table}[ht]
\centering
\caption{Dose delivered to the PTV for different treatments for patient 3.}
\label{tab:patient3-2025}
\begin{tabular}{lllll}
\toprule
Fraction & Overlap & Upper Bound Dose [Gy] & AF Dose [Gy], Cost [ccGy] & Reference Dose [Gy], Cost [ccGy] \\
\midrule
Planning MR & 9.08 cc & --- & --- & --- \\
First Fraction & 19.79 cc & 6 & 6, 0 & 6.6, 11.8 \\
Second Fraction & 6.02 cc & 9 & 9, 18.06 & 6.6, 3.6 \\
Third Fraction & 9.45 cc & 6 & 6, 0 & 6.6, 5.7 \\
Fourth Fraction & 19.59 cc & 6 & 6, 0 & 6.6, 11.8 \\
Fifth Fraction & 12.62 cc & 6 & 6, 0 & 6.6, 7.6 \\
Total Cost & --- & --- & 18 ccGy & 40.5 ccGy \\
\bottomrule
\end{tabular}
\end{table}

\subsection{Patient 49 (40 Gy Prescription)}

Patient 49 was prescribed a total dose of 40 Gy. The observed PTV/OAR overlaps were 4.27 cc in the planning scan and $[5.91, 3.28, 3.87, 4.18, 8.36]$ cc across the five treatment fractions. The allowed dose range spans from 6 Gy to 10 Gy per fraction. Unlike patient 3, where due to the low prescription dose only one high-dose fraction was delivered and the optimal policy was described by a binary decision boundary, patient 49 requires several high-dose fractions.

Figure~\ref{fig:patient49-policy} shows the optimal policy over the five treatment fractions. AF leads to a strategy where two fractions deliver the maximum dose of 10 Gy, two fractions the minimum dose of 6 Gy, and one fraction of 8 Gy is delivered to achieve the 40 Gy prescription dose. Consequently, the optimal policy is no longer characterized by a single threshold, but rather multiple decision boundaries that partition the probability distribution of future overlaps.

In the first fraction, the overlap space is divided into three regions: the algorithm prescribes 10 Gy for overlaps below the 40th percentile, 8 Gy for overlaps between the 40th and 60th percentiles, and 6 Gy otherwise. As the observed overlap on this day was 5.91 cc, above the 60th percentile, the algorithm selects the minimum dose of 6 Gy.

In the second fraction, the policy shows three decision boundaries along the overlap axis, leading to four distinct regions:
\begin{itemize}
\item Below the 26th percentile, a dose of 10 Gy is delivered irrespective of the previously delivered dose.
\item Between the 26th and 50th percentiles, if 10 Gy was delivered in the first fraction, a dose of 8 Gy is delivered; if 8 Gy or less was delivered in the first fraction, overlaps in this range trigger a 10 Gy prescription.
\item Between the 50th and 74th percentiles, if 8 Gy or more was delivered in the first fraction, a dose of 6 Gy is delivered; if 6 Gy was delivered in the first fraction, 8 Gy is prescribed.
\item Above the 74th percentile, a dose of 6 Gy is delivered irrespective of the previously delivered dose.
\end{itemize}
In this fraction, a favorable overlap of 3.28 cc is observed, leading the algorithm to prescribe 10 Gy.

The third fraction exhibits two decision boundaries and three regions. For very low overlap, the algorithm prefers 10 Gy unless the cumulative dose already prevents it. In the intermediate region, the chosen dose depends on the previously accumulated dose, and for overlaps above the 65th percentile the policy generally falls back to 6 Gy, except when a higher residual dose is needed to remain feasible.

By the fourth fraction, only two doses remain to be allocated. Consequently, the policy simplifies to a binary threshold at the 50th percentile: the larger dose level is delivered if the observed overlap is below the mean, and the lower dose level otherwise. Since the observed overlap is 4.18 cc, favorable relative to the updated distribution, the algorithm selects 8 Gy. In the fifth and final fraction, no further adaptation is possible and the residual dose is delivered to match the prescription; for this patient this is 6 Gy.

\begin{figure}[ht]
\centering
\fbox{\parbox{0.8\linewidth}{Figure omitted in this text transcription. See the original PDF for the rendered figure.}}
\caption{Optimal policies and values for patient 49. Panels show the policy for fractions one through five and the expected cost in fraction one. The black crosses mark the observed overlap and the resulting dose decision.}
\label{fig:patient49-policy}
\end{figure}

Table~\ref{tab:patient49-2025} summarizes the AF scheme and the resulting dosimetric cost for patient 49, comparing AF to the upper bound and the uniformly fractionated reference treatment. The AF plan successfully exploits the small overlap volumes of 3.28 cc and 3.87 cc in the second and third fraction by delivering a high dose of 10 Gy. It further avoids delivering high doses during unfavorable configurations. Both the first and last fraction had overlap values above average and received only 6 Gy, minimizing the cost. The 8 Gy fraction is delivered at the intermediate overlap of 4.18 cc in fraction four. Overall, the adaptive strategy achieved a total cost of 37 ccGy, 14.2 ccGy lower than the 51.2 ccGy incurred by the uniform 8 Gy-per-fraction reference treatment. Thereby, AF realizes the optimal treatment given by the upper bound.

\begin{table}[ht]
\centering
\caption{Dose delivered to the PTV for different treatments for patient 49.}
\label{tab:patient49-2025}
\begin{tabular}{lllll}
\toprule
Fraction & Overlap & Upper Bound Dose [Gy] & AF Dose [Gy], Cost [ccGy] & Reference Dose [Gy], Cost [ccGy] \\
\midrule
Planning MR & 4.27 cc & --- & --- & --- \\
First Fraction & 5.91 cc & 6 & 6, 0 & 8, 11.8 \\
Second Fraction & 3.28 cc & 10 & 10, 13.1 & 8, 6.6 \\
Third Fraction & 3.87 cc & 10 & 10, 15.5 & 8, 7.7 \\
Fourth Fraction & 4.18 cc & 8 & 8, 8.4 & 8, 8.4 \\
Fifth Fraction & 8.36 cc & 6 & 6, 0 & 8, 16.7 \\
Total Cost & --- & --- & 37 ccGy & 51.2 ccGy \\
\bottomrule
\end{tabular}
\end{table}

\section{Evaluation Across All Patients}

We applied the AF algorithm to all 58 patients in the dataset. Since the immediate cost is linear in the prescribed dose, AF tends to deliver either the minimum or the maximum dose. An intermediate dose level is delivered only once if needed to match the cumulative prescription. Consequently, for a given cumulative prescription, AF yields the same fractionation scheme and only the order in which the fractions are delivered depends on the observed overlap volumes. For the cumulative prescription levels used in these patients, the algorithm resulted in the following fractionation schemes:

\begin{itemize}
\item For the 33 Gy prescription, the AF plan assigns one 9 Gy fraction and four 6 Gy fractions.
\item For the 35 Gy prescription, the plan includes one 10 Gy, one 7 Gy, and three 6 Gy fractions.
\item The 40 Gy prescription is distributed as two 10 Gy, one 8 Gy, and two 6 Gy fractions.
\item The 45 Gy prescription consists of three 10 Gy, one 9 Gy, and one 6 Gy fraction.
\end{itemize}

Figure~\ref{fig:benefit-distribution} provides an overview of the benefit of AF compared to uniform fractionation across all patients. Most patients show only a small benefit of adaptive fractionation. The mean benefit was 2.9 ccGy. However, selected patients show a substantial benefit. The maximum benefit was achieved for patient 3 (22.5 ccGy). For three patients, AF resulted in worse PTV coverage than uniform fractionation. The largest reduction in PTV coverage was $-4.4$ ccGy. In Figure~\ref{fig:benefit-distribution}b, the benefit is plotted against the standard deviation of the observed overlap volumes. As expected, a clear trend emerges: patients with higher variability in PTV/OAR overlap tend to benefit more from AF.

\begin{figure}[ht]
\centering
\fbox{\parbox{0.8\linewidth}{Figure omitted in this text transcription. See the original PDF for the rendered figure.}}
\caption{Benefit distribution across all patients and correlation between overlap standard deviation, prescription dose, and AF benefit.}
\label{fig:benefit-distribution}
\end{figure}

Table~\ref{tab:best-cases-2025} summarizes the three patients with the largest improvements after patients 3 and 49. For patients 14 and 55 the algorithm correctly identified the lowest overlaps and delivered the high-dose levels, realizing the optimal fractionation scheme. In patient 7 the two fractions with the highest overlap were correctly detected and 6 Gy was prescribed. However, ideally a 10 Gy fraction was delivered in the last fraction rather than the third fraction, which would have resulted in 9.26 ccGy benefit instead of 7.5 ccGy.

\begin{table}[ht]
\centering
\caption{Observed overlaps and delivered doses for patients where AF improved PTV coverage.}
\label{tab:best-cases-2025}
\begin{tabular}{llll}
\toprule
Fraction & Patient 14 & Patient 55 & Patient 7 \\
\midrule
0 & 2.88 cc & 3.00 cc & 5.41 cc \\
1 & 3.35 cc, 6 Gy & 6.41 cc, 6 Gy & 8.20 cc, 6 Gy \\
2 & 3.35 cc, 6 Gy & 4.13 cc, 10 Gy & 9.16 cc, 6 Gy \\
3 & 0.53 cc, 10 Gy & 4.91 cc, 8 Gy & 8.18 cc, 10 Gy \\
4 & 1.92 cc, 7 Gy & 5.66 cc, 6 Gy & 5.45 cc, 10 Gy \\
5 & 3.66 cc, 6 Gy & 3.56 cc, 10 Gy & 7.28 cc, 8 Gy \\
Benefit & 8.8 ccGy & 8.8 ccGy & 7.5 ccGy \\
Explanation & Correctly exploited lowest overlaps in fractions 3 and 4 & Correctly exploited lowest overlaps in fractions 2, 3, and 5 & Correctly identified largest overlaps in fractions 1 and 2, but delivered the high doses in suboptimal order \\
\bottomrule
\end{tabular}
\end{table}

Table~\ref{tab:worst-cases-2025} analyzes the three patients where AF resulted in worse PTV coverage compared to uniform fractionation. Patient 16 had a very small overlap in the planning scan, leading the algorithm to expect similarly favorable geometry in later fractions. As a result, it saved dose for later, but an unusually high overlap in the final fraction led to an overall worse treatment. Patient 5 also had the smallest overlap in the planning scan. The algorithm delayed the delivery of the high-dose fraction until the last fraction, which had a relatively high overlap. Patient 32 presented a very large overlap after the planning scan. The algorithm delivered a large dose in the first and second fraction when the overlap was smaller. However, an even better overlap occurred in the last fraction, which could no longer be exploited. These cases illustrate the limitations of AF, especially when early observations are not representative for the true underlying distribution of overlap volumes.

\begin{table}[ht]
\centering
\caption{Observed overlaps and delivered doses for patients where adaptive fractionation performed worse than uniform fractionation.}
\label{tab:worst-cases-2025}
\begin{tabular}{llll}
\toprule
Fraction & Patient 16 & Patient 5 & Patient 32 \\
\midrule
0 & 1.38 cc & 17.15 cc & 1.56 cc \\
1 & 3.64 cc, 6 Gy & 18.68 cc, 6 Gy & 0.48 cc, 10 Gy \\
2 & 6.73 cc, 6 Gy & 18.35 cc, 6 Gy & 0.22 cc, 10 Gy \\
3 & 3.42 cc, 10 Gy & 28.56 cc, 6 Gy & 0.19 cc, 8 Gy \\
4 & 5.99 cc, 8 Gy & 21.37 cc, 6 Gy & 0.33 cc, 6 Gy \\
5 & 9.17 cc, 10 Gy & 22.98 cc, 9 Gy & 0.00 cc, 6 Gy \\
Benefit & -4.4 ccGy & -3.0 ccGy & -0.7 ccGy \\
Explanation & Small overlap in planning scan not representative for later fractions; high dose had to be delivered in the last fraction despite the largest overlap & Small overlap in planning scan not representative for later fractions & Large overlap in planning scan not representative for later fractions; smallest overlap in last fraction was not exploited \\
\bottomrule
\end{tabular}
\end{table}

In Figure~\ref{fig:upper-bound-gap}, we compare the benefit of AF to the theoretical upper bound that could be achieved if the overlap volumes of all fractions were known in advance. Out of 58 patients, AF achieved the same benefit as the upper bound in 28 patients. The remaining 30 patients showed varying degrees of deviation. Figure~\ref{fig:upper-bound-gap} shows the distribution of differences in benefit between AF and the upper bound for these 30 patients. On average, the benefit of AF is 2.4 ccGy below the upper bound.

Several patients show notable discrepancies. For instance, patient 16 falls short by 22.1 ccGy. The patient shows substantial variation in overlap volumes and, in theory, could have benefited substantially from adaptive fractionation. However, because the small and large overlaps are observed in an unfavorable order, AF actually led to a worse treatment compared to uniform fractionation. Similarly, patient 5 falls short by 13.9 ccGy due to a suboptimal sequence of overlaps. Patient 4 also showed a benefit of approximately 1 ccGy, but the benefit could have been 8.0 ccGy higher if the high-dose fraction had been delivered in the best fraction.

\begin{figure}[ht]
\centering
\fbox{\parbox{0.8\linewidth}{Figure omitted in this text transcription. See the original PDF for the rendered figure.}}
\caption{Histogram of AF benefits relative to the upper bound. Only patients are shown where AF did not deliver the theoretically optimal treatment.}
\label{fig:upper-bound-gap}
\end{figure}

Taken together, these results indicate that selected patients with substantial overlap variation, such as patients 3, 49, 14, 55, and 7, may benefit from AF and the approach can even match the theoretical upper bound on benefit. On the other hand, in a few cases AF may yield worse PTV coverage compared to uniform fractionation. This happens when overlap variations occur in an unfavorable order, especially if the overlap observed in the planning scan is an outlier. For a large group of patients with little anatomical variation, the benefit of AF is modest.

We also investigated the potential benefits of AF in relative numbers. Apart from patients such as 3, 49, and 14, who show large absolute benefits from AF, relevant relative benefits from AF can also be found in patients with smaller PTV and PTV/OAR overlap volumes. For example, patients 20, 24, and 26 all have a lymph node target and a mean PTV volume of approximately 10 cc, an overlap variation of less than 1 cc, and achieve an absolute benefit from AF of 4.05, 3.97, and 5.69 ccGy each. Calculated relative to their mean PTV and PTV/OAR overlap volumes, this represents a sizable benefit in PTV coverage of roughly 2 Gy within the overlap volume, which makes up 20\% of their total PTV.

\section{Dosimetric Evaluation}

The cost defined in equation~\eqref{eq:total-cost} is an approximate measure of PTV underdosage. To assess its limitations and more fully assess the clinical potential of AF, we replanned six patients (3, 14, 18, 25, 26, 49) in the MR-linac treatment planning system, using the dose schedules proposed by the AF algorithm. Our goal was to quantify the true dosimetric benefit, i.e. the additional dose delivered to the PTV compared to the clinically delivered uniformly fractionated treatment. PTV underdosage is now measured by the actual area above the DVH curve for the PTV in each fraction, where the area is integrated up to the prescribed dose in the respective fraction. Summing these values across all fractions provides a total improvement in PTV coverage in units of ccGy per patient.

As an illustrative example, we first revisit patient 3. In Figure~\ref{fig:patient3-dvh}, the DVHs for the first two fractions are shown. The uniformly fractionated plan prescribed 6.6 Gy to the PTV in each fraction. Due to a large OAR/PTV overlap of 19.8 cc in the first fraction (approximately 30\% of the PTV), only around 70\% of the PTV received at least 6 Gy, and about 60\% reached the prescribed dose of 6.6 Gy. This results in an area above the DVH curve of 19.5 ccGy.

In contrast, the AF plan for this first fraction prescribed only 6 Gy. Although the spatial overlap remains the same, AF leads to a reduced area above the DVH of just 6.2 ccGy. The key here is that the evaluation only considers the dose up to the AF-prescribed dose of 6 Gy. In the second fraction, where the overlap was only 6.0 cc, the uniform plan resulted in a smaller area above the DVH of 5.4 ccGy. In the AF plan, a larger dose of 9 Gy was prescribed, and the area above the DVH is 19.2 ccGy. This is still lower than the area above the DVH for the uniformly fractionated plan in the first fraction.

\begin{figure}[ht]
\centering
\fbox{\parbox{0.8\linewidth}{Figure omitted in this text transcription. See the original PDF for the rendered figure.}}
\caption{DVHs of the PTV and GTV for patient 3 under uniformly fractionated and AF plans. Vertical lines indicate prescribed AF doses and uniform doses. Horizontal lines represent the fraction of PTV volume overlapping with the OAR.}
\label{fig:patient3-dvh}
\end{figure}

Figure~\ref{fig:patient3-dose-dist} shows the dose distributions obtained for patient 3, illustrating that a higher dose of 9 Gy can be delivered to most of the PTV in the favorable fraction.

\begin{figure}[ht]
\centering
\fbox{\parbox{0.8\linewidth}{Figure omitted in this text transcription. See the original PDF for the rendered figure.}}
\caption{Replanned dose distributions for patient 3 in a favorable fraction and an unfavorable fraction.}
\label{fig:patient3-dose-dist}
\end{figure}

Table~\ref{tab:patient3-replan} lists the area above the PTV curve for both the uniform and adaptive fractionation plan for all fractions. In the AF plan, only the second fraction receives a high dose of 9 Gy, while all other fractions are prescribed the minimum dose of 6 Gy. Consequently, the AF plan incurs a dose coverage disadvantage in the second fraction only. In all other fractions, AF results in less PTV underdosage compared to uniform fractionation, most notably in the first and fourth fractions, those with the largest PTV/OAR overlap. Overall, adaptive fractionation reduces the cumulative PTV underdosage by 27.6 ccGy compared to uniform fractionation. Compared to the theoretically predicted benefit of 22.4 ccGy, the true benefit even increased.

\begin{table}[ht]
\centering
\caption{Area above the DVH curve for all fractions for patient 3, measured in ccGy. Negative values in the difference column indicate improved coverage with AF.}
\label{tab:patient3-replan}
\begin{tabular}{rrrrr}
\toprule
Fraction & Overlap [cc] & Uniform cost [ccGy] & AF cost [ccGy] & Difference [ccGy] \\
\midrule
1 & 19.97 & 19.5 & 6.2 & -13.3 \\
2 & 6.02 & 5.4 & 19.3 & 13.8 \\
3 & 9.45 & 9.5 & 3.2 & -6.2 \\
4 & 19.59 & 20.0 & 5.6 & -14.4 \\
5 & 12.62 & 13.3 & 5.8 & -7.5 \\
Total & --- & 67.7 & 40.1 & -27.6 \\
\bottomrule
\end{tabular}
\end{table}

As a second example, we examine patient 14, for whom the predicted benefit is 8.8 ccGy. Figure~\ref{fig:patient14-dvh} shows the DVHs for fractions two and three. In fraction two, with an OAR/PTV overlap of 3.35 cc, the uniform plan results in modest PTV underdosage of 5.2 ccGy, while the adaptive plan prescribes a reduced dose of 6 Gy and yields only 0.8 ccGy. In fraction three, which has a minimal overlap of just 0.53 cc, the AF plan delivers a high dose of 10 Gy, with a resulting underdosage of 13.4 ccGy. In comparison, the uniform plan at 7 Gy yields a cost of only 1.0 ccGy.

Although the overlap is significantly lower, the underdosage for a 10 Gy prescription is large in comparison to the underdosage in the uniformly fractionated treatment in other fractions. The reason is illustrated by the corresponding dose distribution: although the volume overlap is small, the bowel wraps around the anterior part of the PTV. Because the dose gradient at the interface of bowel and PTV cannot be arbitrarily steep, parts of the PTV at close distance from the bowel do not receive the prescribed dose of 10 Gy. In the DVH this is seen as falloff of the PTV DVH at large doses, which contributes to the area above the curve.

\begin{figure}[ht]
\centering
\fbox{\parbox{0.8\linewidth}{Figure omitted in this text transcription. See the original PDF for the rendered figure.}}
\caption{DVHs of the PTV and GTV for patient 14 under uniformly fractionated and AF plans.}
\label{fig:patient14-dvh}
\end{figure}

\begin{figure}[ht]
\centering
\fbox{\parbox{0.8\linewidth}{Figure omitted in this text transcription. See the original PDF for the rendered figure.}}
\caption{Dose distributions for patient 14 in fraction 3, comparing the replanned AF dose of 10 Gy to the original 7 Gy plan.}
\label{fig:patient14-dose-dist}
\end{figure}

Table~\ref{tab:patient14-replan} summarizes the per-fraction PTV underdosage for both plans. While the AF plan achieves better coverage in three of the five fractions, the 10 Gy prescription in fraction 3 yields a disproportionately high cost. As a result, the total underdosage for AF (19.0 ccGy) is nearly identical to that of the uniform plan (18.9 ccGy). The predicted benefit of 8.8 ccGy could not be realized in this patient.

\begin{table}[ht]
\centering
\caption{Area above the DVH curve per fraction for patient 14, measured in ccGy. Negative values in the difference column indicate a benefit of AF.}
\label{tab:patient14-replan}
\begin{tabular}{rrrrr}
\toprule
Fraction & Overlap [cc] & Uniform cost [ccGy] & AF cost [ccGy] & Difference [ccGy] \\
\midrule
1 & 3.35 & 4.5 & 1.0 & -3.5 \\
2 & 3.35 & 5.2 & 0.8 & -4.4 \\
3 & 0.53 & 1.0 & 13.4 & 12.4 \\
4 & 1.92 & 2.8 & 2.8 & 0.0 \\
5 & 3.66 & 5.4 & 1.0 & -4.4 \\
Total & --- & 18.9 & 19.0 & 0.1 \\
\bottomrule
\end{tabular}
\end{table}

Table~\ref{tab:all-replanned-2025} summarizes the results for all six replanned patients. For most patients, the area above the DVH is smaller for the AF plan compared to the clinically delivered uniformly fractionated plans, showing that the benefit of AF can be realized. However, the benefit tends to be smaller than what is predicted by the algorithm.

\begin{table}[ht]
\centering
\caption{Area above the DVH for all six replanned patients for the uniformly fractionated plan and the AF plan. Negative benefit values indicate reduced underdosage with AF.}
\label{tab:all-replanned-2025}
\begin{tabular}{rrrrrr}
\toprule
Patient & Prescription [Gy] & Uniform [ccGy] & AF [ccGy] & Actual Benefit [ccGy] & Predicted Benefit [ccGy] \\
\midrule
3 & 33 & 67.7 & 40.1 & -27.6 & -22.4 \\
14 & 35 & 18.9 & 19.0 & 0.1 & -8.8 \\
18 & 35 & 6.2 & 2.1 & -4.1 & -3.6 \\
25 & 35 & 1.4 & 0.1 & -1.3 & -1.6 \\
26 & 35 & 13.6 & 9.5 & -4.1 & -5.7 \\
49 & 40 & 79.9 & 71.6 & -8.3 & -14.2 \\
\bottomrule
\end{tabular}
\end{table}

\section{Discussion}

\subsection{Summary of Main Findings}

We consider challenging cases for MR-guided SBRT where the PTV prescription dose is higher than the dose constraint for an OAR that overlaps with the PTV, requiring compromises in PTV coverage. We introduced an approach to AF that exploits interfraction variation in the overlap volume between the PTV and the dose-limiting OAR, aiming to minimize underdosage within the PTV. This is achieved by delivering higher doses in fractions with favorable PTV/OAR geometries, i.e. small overlap volumes, and lower doses on days with unfavorable geometries, i.e. large overlap volumes. An MDP-based algorithm was introduced to determine the daily delivered dose.

Across all 58 patients that were evaluated, AF provided a dosimetric benefit in a total of 55 cases that showed reduced PTV underdosage compared to uniform fractionation. The mean benefit for all patients was 2.9 ccGy, of which a small subset achieved more substantial improvements, up to 22.5 ccGy. The magnitude of the observed benefits was strongly correlated with the degree of variation in PTV/OAR overlap, illustrating the relevance of this parameter as a dominant predictor of AF performance. Conversely, patients with low or unfavorably ordered overlap variations exhibited limited benefit or even a dosimetric disadvantage compared to uniform fractionation. Overall, these findings indicate that overlap-guided AF can achieve moderate dosimetric improvements for most patients and larger benefits for selected cases with pronounced anatomical variability.

\subsection{Relation to Other Works}

Previous works introduced the concept of AF \cite{ref10,ref11,ref12,ref13}. These methods have in common that the fraction doses are adapted with the goal of exploiting daily anatomical variations. However, the approach presented here differs regarding its goal and the situation in which it is applicable. Prior works quantified daily variations in terms of the sparing factor, defined as the ratio of OAR BED and PTV BED, and the goal was to maximize the cumulative PTV BED \cite{ref13}. However, for the sparing factor to vary between fractions, the distance between PTV and OAR must change. In situations where PTV and OAR overlap in all fractions, the sparing factor is approximately constant. Thus, these earlier approaches to AF were not applicable in these situations. The variant of AF introduced here addresses this shortcoming.

Other research groups have also explored ideas to adapt fraction doses with different objectives. In a recent study by van Lieshout et al., the primary goal was optimizing treatment efficiency in the context of online adaptive radiotherapy by adapting the number of delivered fractions \cite{ref15}. In an in-silico cohort of patients receiving SBRT for abdominopelvic lymph node oligometastases, patients showing advantageous PTV/OAR geometry at treatment were replanned using higher fraction doses with the goal of shortening overall treatment duration while respecting OAR dose constraints. According to the authors, this concept will now be carried forward into a clinical trial.

\subsection{Limitations and Future Work}

The area above the DVH curve is used here as a measure of PTV underdosage. In the AF algorithm, this is approximated by a square as illustrated in Figure~\ref{fig:dvh-schematic}. As described in section~6, the AF algorithm may underestimate the true underdosage as measured by the area above the DVH that is obtained after adaptive replanning with the prescription doses obtained from the algorithm. The reason is that the AF algorithm does not correctly account for the underdosage that occurs in the non-overlapping part of the PTV due to the finite steepness of the dose gradient that can be achieved at the interface to the OAR. The magnitude of this additional underdosage depends on the fraction dose but also on the geometry of the target volume in relation to the dose-limiting OAR. Future work may develop improved methods to estimate the area above the DVH, taking into account the falloff of the DVH at large doses.

A practical difficulty regarding a future clinical implementation of AF in the form of a pilot phase I trial concerns the observation that only relatively few patients have a substantial dosimetric benefit from AF that is potentially clinically relevant. For most patients, the daily variation in the overlap volume turns out to be too small. The ability to predict the amount of overlap variation and the benefit of AF for a given patient a priori would therefore be helpful. For the patient cohort analyzed in this work, no clear relationship between overlap variation and PTV volume or tumor site was apparent. Future work may aim to identify surrogate parameters or indicators that can more reliably estimate AF benefit based on pre-treatment imaging.

\section{Conclusion}

Overlap-variation-guided AF is an approach for optimizing PTV coverage by adapting dose delivery to daily variations in PTV/OAR geometry in patients receiving MR-guided SBRT for abdominal and pelvic tumors with large interfractional motion that overlap with dose-limiting OARs. By assigning higher doses to fractions with smaller overlap volumes and lower doses to those with larger ones, the method improves target coverage without compromising OAR sparing. While this approach yields only moderate dosimetric advantages in most patients, some cases with pronounced anatomical variability show substantial benefit.

\appendix
\section{Appendix}

The full implementation of the adaptive fractionation algorithm, together with a user interface, all analysis scripts, and anonymized data used in this study, is available on GitHub:
\url{https://github.com/YoelPH/Overlap_adaptfx/tree/paper-version}

\begin{figure}[ht]
\centering
\fbox{\parbox{0.8\linewidth}{Figure omitted in this text transcription. See the original PDF for the rendered figure.}}
\caption{Mean PTV volumes in the patient cohort. The colors represent the respective tumor locations.}
\label{fig:appendix-ptv-volumes}
\end{figure}

\begin{figure}[ht]
\centering
\fbox{\parbox{0.8\linewidth}{Figure omitted in this text transcription. See the original PDF for the rendered figure.}}
\caption{Variation in overlap volumes for all patients. Each star represents the total calculated overlap volume for one fraction, including the initial simulation. The red line represents the mean overlap volume throughout each treatment course.}
\label{fig:appendix-all-overlaps}
\end{figure}

\subsection*{Appendix Table Note}

The original manuscript contains a large patient-by-patient appendix table (Table 8) listing AF plans, tumor sites, upper bounds, overlaps, overlapping OARs, and benefits for all 58 patients. That source data remains available in:
\begin{itemize}
\item \texttt{docs/papers/2025/Overlap\_Adaptive\_Fractionation\_2025Paper.pdf}
\item \texttt{docs/papers/2025/Overlap\_Adaptive\_Fractionation\_2025Paper.md}
\end{itemize}
It has not been fully re-typeset in this first-pass LaTeX transcription.

\begin{thebibliography}{99}

\bibitem{ref1} Acharya S et al. Online magnetic resonance image guided adaptive radiation therapy: First clinical applications. \textit{International Journal of Radiation Oncology Biology Physics} 94, 394--403 (2016).
\bibitem{ref2} Mutic S and Dempsey J F. The ViewRay system: magnetic resonance-guided and controlled radiotherapy. \textit{Seminars in Radiation Oncology} 24, 196--199 (2014).
\bibitem{ref3} Raaymakers B W et al. First patients treated with a 1.5 T MRI-linac: clinical proof of concept of a high-precision, high-field MRI guided radiotherapy treatment. \textit{Physics in Medicine \& Biology} 62, L41 (2017).
\bibitem{ref4} Kl\"uter S. Technical design and concept of a 0.35 T MR-linac. \textit{Clinical and Translational Radiation Oncology} 18, 98--101 (2019).
\bibitem{ref5} Palacios M A et al. Role of daily plan adaptation in MR-guided stereotactic ablative radiation therapy for adrenal metastases. \textit{International Journal of Radiation Oncology Biology Physics} 102, 426--433 (2018).
\bibitem{ref6} Mayinger M et al. Benefit of replanning in MR-guided online adaptive radiation therapy in the treatment of liver metastasis. \textit{Radiation Oncology} 16 (2021).
\bibitem{ref7} Pavic M et al. MR-guided adaptive stereotactic body radiotherapy of primary tumor for pain control in metastatic pancreatic ductal adenocarcinoma: an open randomized, multicentric, parallel group clinical trial (MASPAC). \textit{Radiation Oncology} 17 (2022).
\bibitem{ref8} Chuong M D et al. Stereotactic MR-guided on-table adaptive radiation therapy for borderline resectable and locally advanced pancreatic cancer: a multi-center, open-label phase 2 study. \textit{Radiotherapy and Oncology} 191, 110064 (2024).
\bibitem{ref9} Henke L et al. Phase I trial of stereotactic MR-guided online adaptive radiation therapy for the treatment of oligometastatic or unresectable primary malignancies of the abdomen. \textit{Radiotherapy and Oncology} 126, 519--526 (2018).
\bibitem{ref10} Lu W, Chen M, Chen Q, Ruchala K and Olivera G. Adaptive fractionation therapy: I. Basic concept and strategy. \textit{Physics in Medicine \& Biology} 53, 5495 (2008).
\bibitem{ref11} Chen M, Lu W, Chen Q, Ruchala K and Olivera G. Adaptive fractionation therapy: II. Biological effective dose. \textit{Physics in Medicine \& Biology} 53, 5513 (2008).
\bibitem{ref12} Ramakrishnan J, Craft D, Bortfeld T and Tsitsiklis J N. A dynamic programming approach to adaptive fractionation. \textit{Physics in Medicine \& Biology} 57, 1203 (2012).
\bibitem{ref13} Haas Y P, Ludwig R, Bello R D, Tanadini-Lang S and Unkelbach J. Adaptive fractionation at the MR-linac. \textit{Physics in Medicine \& Biology} 68, 035003 (2023).
\bibitem{ref14} Sutton R S and Barto A G. \textit{Reinforcement Learning: An Introduction}. MIT Press, Cambridge, MA, USA (2018).
\bibitem{ref15} van Lieshout E et al. Reducing number of treatment fractions for patients with abdominal lymph node oligometastases: the need for online adaptive radiation therapy to provide personalized adaptive fractionation. \textit{International Journal of Radiation Oncology Biology Physics} 122, 721--728 (2025).

\end{thebibliography}

\end{document}
